\documentclass[11pt]{article}
\usepackage[margin=1in]{geometry}
\usepackage{amsmath,amssymb,amsthm}
\usepackage{hyperref}
\newtheorem{theorem}{Theorem}
\newtheorem{definition}{Definition}
\title{Quantile Index Preservation under Bounded Score Perturbations}
\author{Research Architecture Runtime \\ \texttt{operator/quantile}}
\date{2026-07-25}
\begin{document}
\maketitle
\begin{abstract}
We study $\tau$-quantile selection as the unique strict $k$-th order statistic on a
finite real score vector. If every pairwise gap exceeds $2\varepsilon$, the selected
index is invariant under $\|\delta\|_\infty\le\varepsilon$. Lean \texttt{LEAN\_FULL}
on Mathlib $\mathbb{R}$ (\texttt{REAL\_MATHLIB}).
\end{abstract}
\section{Operator}
\begin{definition}
For $s\in\mathbb{R}^n$ ($n\ge 2$), index $i$ is the unique strict $k$-th smallest when
exactly $k$ scores are strictly below $s_i$ and no other index shares $s_i$.
A $\tau$-quantile uses $k=\lfloor\tau(n-1)\rfloor$.
\end{definition}
\section{Theorems}
\begin{theorem}[Preservation]
If $i$ is a unique strict $k$-th smallest and all pairwise gaps exceed $2\varepsilon$,
then every $\|\delta\|_\infty\le\varepsilon$ perturbation preserves that status.
\end{theorem}
\begin{theorem}[Sharpness]
If some rival $j\neq i$ has $|s_i-s_j|\le 2\varepsilon$, an admissible $\delta$
destroys unique strict $k$-th-smallest status of $i$ (tie adversary).
\end{theorem}
\section{Proofs}
Pairwise order vs $i$ is preserved under the gap hypothesis; \texttt{countLT} congruence
and uniqueness follow. Sharpness forces a value-tie between $i$ and a close rival.
\section{Comparison}
Relative to Argmax margin: same $2\varepsilon$ gap cost, but ranking/count geometry
rather than unique-max. Relative to Threshold: multi-cut ordering instead of a single cut.
\section{Lean / certificate}
\texttt{Research.Operators.Quantile.Preservation} reducing to
\texttt{OrderStat.KthMargin}; certificate
\texttt{lean/certificates/quantile/quantile-margin/} (\texttt{LEAN\_FULL}; Mathlib $\mathbb{R}$).
\section{Limitations}
Mathlib $\mathbb{R}$ (\texttt{REAL\_MATHLIB}); unique-value hypothesis; no privacy claims.
\end{document}
